% Abhakseite „Das kann ich“ – Zone nur im Gesamt (\mitzone)
%% TODO Ich-kann-Sätze: Platzhalter wie in den Aufgabendateien
\begin{abhakseite}
\ifdefined\mitzone
\abhakgruppe{Kennst du schon}
\abhak{1}{Vierfeldertafel füllen (Ränder, Felder, bedingte Angaben)}
\abhak{2}{Bedingte Wahrscheinlichkeit als Quotient bilden}
\abhak{3}{Anteile aus absoluten Häufigkeiten bilden und kürzen}
\abhak{4}{Pfadterme mit Parameter im Baumdiagramm (Ast als Faktor, Pfad als Produkt, Randwahrscheinlichkeit als Pfadsumme)}
\abhak{5}{Lineare und quadratische Gleichungen aufstellen und lösen, Lösungen sieben (Nulllösung verwerfen)}
\abhak{6}{Pfadregeln und Gegenereignis mit und ohne Zurücklegen}
\abhak{7}{Vierfeldertafel füllen (Ränder, Felder, bedingte Angaben) – Fehler finden}
\abhak{8}{Vierfeldertafel füllen (Ränder, Felder, bedingte Angaben) – selbst rechnen}
\fi
\abhakgruppe{Einheit 1 · Unabhängigkeit prüfen: die Produktregel P(A ∩ B) = P(A) · P(B) an Anteilen oder absoluten Häufigkeiten prüfen (Tafel oder Anzahlen erst beschaffen}
\abhak{9}{„Prüfen oder nutzen?“}
\abhak{10}{Unabhängigkeit prüfen}
\abhak{11}{Unabhängigkeit prüfen – Fehler finden, Begründen, Anwendung, Darstellungswechsel}
\abhakgruppe{Einheit 2 · Rückwärts}
\abhak{12}{Rückwärts}
\abhak{13}{Rückwärts – Fehler finden, Begründen, Anwendung}
\abhakgruppe{Einheit 3 · Unabhängigkeit als Argument: die Ausgleichs-Fehlvorstellung widerlegen}
\abhak{14}{Als Argument}
\abhak{15}{Als Argument – Fehler finden, Begründen}
\end{abhakseite}
